\chapter{Preliminaries}

\begin{figure}
\centering
\includegraphics[width=0.39\textwidth]{figures/example_boards/preliminary.pdf}
\caption[Example instance of the Polyomino Assembly Problem]{Annotated example instance of the Polyomino Assembly Problem. The blocked positions on the board are shown as grey squares and the tiles as colored squares. The three connected tiles form a non-trivial polyomino.}
\label{fig:example_instance}
\end{figure}

\paragraph{Board} A \emph{board} $B$ is a rectangular region of $\mathbb{Z}^2$ where each position is either marked as \emph{open} or \emph{blocked}. Blocked positions represent obstacles and cannot contain tiles. Furthermore, the induced grid graph $G_{B}$ is defined as the graph in which the open positions of $B$ are nodes and two nodes are connected if and only if their distance is 1. $B$ is called \emph{connected} if and only if $G_{B}$ is connected.

\paragraph{Tiles and Glues} A tile $t= (p, g)$ is a unit square centered on $p$ with edge glues $g$, where $p$ is an open position on the board and $g$ is a quadruple $(g_{N}(t), g_{E}(t), g_{S}(t), g_{W}(t))$ of glues. $g_{d}(t)$ denotes the glue on the edge of $t$ facing in the cardinal direction $d \in \{N,E,S,W\}$. Each glue is an element of a finite alphabet $\Sigma$, which also contains a special \texttt{null} glue. A glue function $G \colon \Sigma \to \{0, 1\}$ defines which glues stick to each other. Two glues $g_{1}, g_{2} \in \Sigma$ stick to each other with respect to $G$ if and only if $G(g_{1}, g_{2}) = 1$. Furthermore, $G$ has the properties $G(g_{1}, g_{2}) = G(g_{2}, g_{1})$ and $G({g_{1}, \texttt{null}}) = 0$ for all $g_{1}, g_{2} \in \Sigma$. Tiles \emph{bond}, if and only if they are located on adjacent positions and the glues on the shared edge stick together. The bond is considered permanent, and the tiles hereafter move as a unit.

\paragraph{Polyomino} A \emph{polyomino} $P$ is a finite set of tiles that forms a connected component in the graph in which bonded tiles are adjacent. A polyomino containing exactly one tile is called trivial. The position of $P$ is defined as the position of the top-leftmost tile in $P$.

\paragraph{Configuration} A \emph{configuration} is a tuple $C = (B, T)$, consisting of a a board $B$ and a set of tiles $T$. The positions of tiles must be non-overlapping and open in $B$.

\paragraph{Step} A \emph{step} is a transformation between two configurations that moves every tile by one unit distance in one of the cardinal directions unless that tile is part of a blocked polyomino.
Formally, a polyomino is \emph{blocked} in the direction $d$ if and only if it is not contained in the maximal subset of polyominoes $P$ such that moving all tiles in $\bigcup P$ by one unit distance in the direction $d$ would neither cause tiles to overlap with each other nor with blocked positions. For simplicity, a step can be denoted by the direction $d \in \{N, E, S, W\}$ in which the polyominoes are moved. A \emph{step sequence} is a series of steps.

\section{Problem Definitions}

\paragraph{Polyomino Assembly Problem}
Given a configuration $C= (B, T)$ and a set of open positions $X \subseteq B$ that is connected in $G_{B}$ the Polyomino Assembly Problem asks for either a step sequence $S$ such that consecutively applying the steps of $S$ to $C$, results in a configuration in which a polyomino $P$ exists that satisfies $X = \{p_{1}, p_{2},\ldots \}$, where $p_{1}, p_{2},\ldots$ are the positions of the tiles in $P$, or a proof that the configuration is unreachable.

\paragraph{Fixed Seed Tile Polyomino Assembly Problem}
Given a configuration $C= (B, T)$, a set of open positions $X \subseteq B$ that is connected in $G_{B}$ and a fixed tile $t_{\text{fixed}} \in T$ the Polyomino Fixed Seed Tile Polyomino Assembly Problem is defined analogously to the Polyomino Assembly Problem, except that $t_{\text{fixed}}$ does not move under step transformation and tiles can only bond if the resulting polyomino contains $t_{\text{fixed}}$.


\section{Distance Definitions}
As a basis for the definition of our heuristic functions, we use the following definitions of distance between positions, tiles and the target shape.
Given an instance of the Polyomino Assembly Problem with a configuration $C= (B, T)$, and target shape $X$. Let $p = (p_1, p_2), q = (q_1, q_2)$ be open positions in B and $t=(p_t, g_t) \in T$.
\begin{itemize}
	\setlength\itemsep{0em}
	\item $d(p, q)$ is the length of a shortest path from $p$ to $q$ in $G_{B}$.
	\item $d(p)$ is the length of a shortest path in $G_{B}$ from $p \in B$ to the closest node in $X$.
	\item $d_{1}(p, q) \coloneqq |p_1 - q_1| + |p_2 - q_2|$ is called the taxicab distance.
\end{itemize}
For the sake of brevity, we define $d(t) \coloneqq d(p_t)$ and analogously for the other distance definitions.